\chapter{Work Log}

This appendix records the engineering workflow used during the early project stages. It is intended as a traceability aid during development and may be shortened or removed in the final submitted thesis.

\section*{Stage 1: Project Skeleton}

The initial directory structure was created, including folders for aircraft data, plant dynamics, controller design, guidance logic, validation scripts, feature tests, tooling, plots, Simulink models and documentation. A LaTeX thesis draft was added so that documentation could evolve together with the implementation.

\section*{Stage 2: Modelling Conventions}

The main reference frames, units, state vectors, input vectors and control-surface sign conventions were defined. The full plant state was extended to include actual engine thrust and actuator states.

\section*{Stage 3: Aircraft Data Layer}

The first reusable aircraft data modules were implemented. Geometry, mass, inertia, standard atmosphere, reference cases and a simple aerodynamic coefficient model were introduced. Feature checks were updated to verify the consistency of these modules.

\section*{Stage 4: Staged Aerodynamic Database}

The aerodynamic model was reorganised around a database-style interface. The initial analytical model remains available as a low-confidence seed, while the new database layer records the intended source hierarchy: public B777-like geometry and mass data, NASA CRM clean transonic aerodynamics, DATCOM derivative estimates, NASA CRM-HL high-lift data and selected OpenFOAM verification cases. The first NASA CRM CFD table at $M=0.85$ was added as a loader cross-check, and the active clean-cruise baseline was then upgraded to a NASA CRM NTF197/TWICS Mach-alpha grid for \texttt{CL}, \texttt{CD} and \texttt{Cm}. Dedicated CSV readers were introduced so that NASA CRM force and moment data can be added as source files rather than manually copied coefficient vectors. Seed derivative and high-lift tables were separated from the main database so that DATCOM and CRM-HL data can later replace them without changing the plant interface. The derivative layer was then upgraded from passive metadata to an active DATCOM-ready residual table used for sideslip, rate and control effects inside the CRM clean-grid domain. A reproducible Digital DATCOM workflow was then added: the public PDAS archive is downloaded and compiled locally, committed B777-like clean input decks are run across a five-point Mach grid, and each output is parsed automatically into a candidate derivative table with a separate extraction report documenting selected rows, skipped blank fields and \texttt{NDM} entries. The per-Mach derivative files were then consolidated into a Mach-dependent candidate table with a grid-level source identifier, retained per-run source identifiers and a no-extrapolation interpolation policy. Promotion gates were then added so that source traceability and sign checks can pass while activation remains blocked until active-interface coverage, control derivatives and validation evidence are available. The DATCOM workflow was then extended with generated elevator and aileron control-effectiveness decks. The new control-grid runner executes \texttt{SYMFLP} and \texttt{ASYFLP} cases over the same Mach stations, extracts elevator and aileron candidates, consolidates them into a separate control-grid table and records the entries that are not direct DATCOM printouts. A derived gap-fill builder was then added for $C_{D_\beta}$, $C_{Y_r}$, rudder derivatives and lateral control drag, using the DATCOM candidate tables and B777-like geometry so that active-interface coverage can be checked without silently using seed values. A clean-cruise longitudinal trim-validation helper was then added for the DATCOM candidate path; three representative Mach cases pass the current force, moment and static-equivalent thrust checks. A restricted longitudinal modal-validation helper was then added for the same cases; each case produces stable phugoid and short-period mode pairs, while the DATCOM warning review remains open. The high-lift layer was then upgraded to a CRM-HL/DATCOM-ready source table containing clean, approach and landing configurations with flap, slat, gear, increment and \texttt{CLmax} metadata. The plant-facing aerodynamic load adapter was added to compute dynamic pressure, body-axis forces, aerodynamic moments and centre-of-gravity moment shifts. The aerodynamic feature test now verifies the database references, source inventory, derivative inventory, CRM grid values, DATCOM-ready residual effects, DATCOM candidate-run artifacts, trim and modal validation evidence, CRM-HL/DATCOM-ready high-lift behaviour, force signs and CG-shift logic.
